% SEGMENT OLP-0541-S01
% Part: set-theory
% Chapter: z
% Section: pairs

\documentclass[../../../include/open-logic-section]{subfiles}

\begin{document}

\olfileid{sth}{z}{pairs}
\olsection{ஜோடிகள்}

அடுத்து கருத வேண்டிய அடிகோள் பின்வருவது:

\begin{axiom}[ஜோடிகள்]
எந்தக் கணங்கள் $a, b$ க்கும், $\{a, b\}$ என்ற கணம் உள்ளது.
\[
	\forall a \forall b \exists P \forall x (x \in P \liff (x = a \lor x = b))
\]
\end{axiom}

படிநிலைமுறைக் கருத்தாக்கத்தைப் பயன்படுத்தி இந்த அடிகோளை எவ்வாறு
நியாயப்படுத்துவது என்பது இதோ. படிநிலை $S$ இல் $a$ கிடைக்கிறது என்றும்,
படிநிலை $T$ இல் $b$ கிடைக்கிறது என்றும் கொள்க. $S$, $T$ ஆகியவற்றுள்
பிந்திய படிநிலையை $M$ என்க. அப்போது $a$, $b$ ஆகிய இரண்டும் படிநிலை $M$
இல் கிடைப்பதால், $\{a,b\}$ என்ற தொகுப்பு $M$ க்குப் பிந்திய எந்தப்
படிநிலையிலும் உருவாகக்கூடிய ஒரு கணமாகும் (அவற்றுள் பிந்தியதே அது).

ஆனால் சற்றுப் பொறுங்கள்!{} $M$ க்குப் பிந்திய படிநிலைகள் எவையும்
\emph{உள்ளன} என்று ஏன் எடுகோள வேண்டும்? அவை இல்லை என்றால் நமது
நியாயப்படுத்தல் தோல்வியடையும். ஆகவே ஜோடிகள் அடிகோளை நியாயப்படுத்த,
\olref[sth][z][story]{sec} இல் கூறிய கதையில் மேலும் ஒரு கொள்கையைச்
சேர்க்க வேண்டும்; அது:
\begin{enumerate}
	\item[] \stagessucc. கடைசிப் படிநிலை இல்லை.
\end{enumerate}
இந்தக் கொள்கை நியாயப்படுத்தப்பட்டதா? கடைசிப் படிநிலை இல்லை என்று
ஷோன்ஃபீல்டின் கதை \emph{வெளிப்படையாகக்} கூறவில்லை. இருந்தாலும், அது
(கறாராகச் சொன்னால்) நமது கதைக்கான ஒரு கூடுதல் சேர்க்கையாக
இருந்தாலும், கணங்கள் படிநிலைகளில் உருவாகின்றன என்ற அடிப்படைக் கருத்துடன்
நன்கு பொருந்துகிறது. தொடர்ந்து வருவதில் இதனை வெறுமனே ஏற்றுக்கொள்வோம்.
எனவே ஜோடிகள் அடிகோளையும் ஏற்றுக்கொள்வோம்.
% SOURCE CORRECTION TA-STH-007: restore Schoenfield's name where the source misspells it “Shoenfield”.

இந்தப் புதிய அடிகோளைக் கொண்டு மேலும் பல கணங்கள் உள்ளன என நிறுவலாம்.
எடுத்துக்காட்டாக:

\begin{prop}\ollabel{prop:pairsconsequences}
எந்தக் கணங்கள் $a$, $b$ க்கும், பின்வரும் கணங்கள் உள்ளன:
	\begin{enumerate}
		\item\ollabel{singleton} $\{a\}$
		\item\ollabel{binunion} $a \cup b$
		\item\ollabel{tuples} $\tuple{a, b}$
	\end{enumerate}
\end{prop}

\begin{proof}
\olref{singleton}. ஜோடிகள் அடிகோளின்படி $\{a, a\}$ உள்ளது; உறுப்புசார்
சமத்துவத்தின்படி இதுவே $\{a\}$.

\olref{binunion}. ஜோடிகள் அடிகோளின்படி $\{a, b\}$ உள்ளது. இப்போது
$a \cup b = \bigcup \{a, b\}$ என்பது சேர்ப்பு அடிகோளின்படி உள்ளது.

\olref{tuples}. \olref{singleton} இன்படி $\{a\}$ உள்ளது. ஜோடிகள்
அடிகோளின்படி $\{a, b\}$ உள்ளது. இப்போது மீண்டும் ஜோடிகள் அடிகோளின்படி
$\{\{a\}, \{a, b\}\} = \tuple{a, b}$ உள்ளது.
\end{proof}

\begin{prob}
எந்தக் கணங்கள் $a, b, c$ க்கும் $\{a, b, c\}$ என்ற கணம் உள்ளது எனக்
காட்டுக.
\end{prob}

\begin{prob}
எந்தக் கணங்கள் $a_1, \ldots, a_n$ க்கும் $\{a_1, \ldots,
a_n\}$ என்ற கணம் உள்ளது எனக் காட்டுக.
\end{prob}

%\begin{proof} ஜோடிகள் அடிகோளை இருமுறை பயன்படுத்தி, $\{\{a_1, a_2\}, \{a_1,
%   a_3\}\}$ உள்ளது. சேர்ப்பு மற்றும் உறுப்புசார் சமத்துவத்தின்படி, $\bigcup \{\{a_1,
%   a_2\}, \{a_1, a_3\}\} = \{a_1, a_2, a_3\}$ உள்ளது. தேவையான அளவு
%   முறை இந்த உத்தியைத் திரும்பப் பயன்படுத்துக. \end{proof}

\end{document}
